\documentclass[11pt,letterpaper]{article}
\usepackage[margin=0.75in]{geometry}
\usepackage[T1]{fontenc}
\usepackage{textcomp,lmodern,microtype,xcolor,enumitem,listings,fancyhdr,hyperref}
\definecolor{accent}{HTML}{184B59}
\definecolor{light}{HTML}{F1F5F6}
\hypersetup{colorlinks=true,urlcolor=accent,pdftitle={Quick Python and VS Code Setup}}
\setlength{\parindent}{0pt}
\setlength{\parskip}{6pt}
\setlist{leftmargin=*,itemsep=4pt}
\lstset{basicstyle=\ttfamily\small,backgroundcolor=\color{light},frame=single,rulecolor=\color{light},framesep=6pt,breaklines=true,columns=fullflexible,keepspaces=true,showstringspaces=false,upquote=true}
\pagestyle{fancy}\fancyhf{}
\fancyhead[L]{\small Student software setup}
\fancyfoot[R]{\small\thepage}
\renewcommand{\headrulewidth}{0pt}
\setlength{\headheight}{14pt}
\newcommand{\heading}[2]{{\small\color{accent}\textbf{#1}}\par{\LARGE\bfseries #2}\par}
\begin{document}
\sloppy
\heading{WINDOWS}{Quick Python and VS Code setup}
Follow this page on Windows, or page 2 on Mac. Then finish on page 3.
Allow about 20--30 minutes, plus download time.

\section*{1 \quad Install Python from the Microsoft Store}
Open the \textbf{Microsoft Store}. Search for \textbf{Python Install Manager}
and install the one published by the \textbf{Python Software Foundation}.
Open \textbf{Windows PowerShell} from the Start menu. Copy and paste this
block, press Enter, and wait for it to finish:
\begin{lstlisting}
pymanager install 3.13
py -3.13 --version
\end{lstlisting}
You should see \texttt{Python 3.13.x}. Close PowerShell when finished.

\section*{2 \quad Install VS Code}
Download and install \href{https://code.visualstudio.com/Download}{Visual Studio Code}.
Open it, click \textbf{Extensions} in the left sidebar, and install these two
extensions by \textbf{Microsoft}: \textbf{Python} and \textbf{Jupyter}.

\section*{3 \quad Create a project folder}
Create a folder named \texttt{RTS-Research} in Documents.
In VS Code, choose \textbf{File > Open Folder} and open that folder.
Choose \textbf{Terminal > New Terminal}. Select \textbf{PowerShell}
from the terminal dropdown if needed.

\section*{4 \quad Install all libraries at once}
Copy and paste this whole block into that PowerShell terminal. Press Enter
if the last line is waiting. Leave it open and wait for installation to finish:
\begin{lstlisting}
py -3.13 -m venv .venv
.\.venv\Scripts\python.exe -m pip install --upgrade pip
.\.venv\Scripts\python.exe -m pip install numpy pandas scipy `
  matplotlib seaborn scikit-learn openpyxl ipykernel notebook
.\.venv\Scripts\python.exe -m pip check
\end{lstlisting}
\textbf{Done when:} the last line says \texttt{No broken requirements found}.
This installs tools for data cleaning, Excel, plotting, ML, and Jupyter.
The \texttt{.venv} folder keeps your project libraries together; no activation is needed.

\textbf{Next: skip the Mac page and go to page 3.}

\newpage
\heading{MAC}{Quick Python and VS Code setup}
Windows users: skip this page.

\section*{1 \quad Install Python}
Open the \href{https://www.python.org/downloads/release/python-31316/}{Python 3.13.16 download page}.
Under \textbf{Files}, download the \textbf{macOS installer}.
Open the downloaded \texttt{.pkg} file and follow the installation steps.

In Finder, open \textbf{Applications > Python 3.13} and double-click
\textbf{Install Certificates.command}. Let it finish.

\section*{2 \quad Install VS Code}
Download \href{https://code.visualstudio.com/Download}{Visual Studio Code for Mac}.
Unzip it if needed, drag the app into \textbf{Applications}, and open it.
Click \textbf{Extensions} in the sidebar and install these two extensions
by \textbf{Microsoft}: \textbf{Python} and \textbf{Jupyter}.

\section*{3 \quad Create a project folder}
Create a folder named \texttt{RTS-Research} in Documents.
In VS Code, choose \textbf{File > Open Folder} and open that folder.
Then choose \textbf{Terminal > New Terminal}.

\section*{4 \quad Install all libraries at once}
Copy and paste this whole block into the terminal. Press Enter if the last
line is waiting. Leave it open and wait for installation to finish:
\begin{lstlisting}
python3.13 -m venv .venv
./.venv/bin/python -m pip install --upgrade pip
./.venv/bin/python -m pip install numpy pandas scipy \
  matplotlib seaborn scikit-learn openpyxl ipykernel notebook
./.venv/bin/python -m pip check
\end{lstlisting}
\textbf{Done when:} the last line says \texttt{No broken requirements found}.
This installs tools for data cleaning, Excel, plotting, ML, and Jupyter.
The \texttt{.venv} folder keeps your project libraries together; no activation is needed.

\textbf{Next: go to page 3.}

\newpage
\heading{EVERYONE}{Open a notebook and check your setup}
\section*{5 \quad Create your first notebook}
\begin{enumerate}
\item Open the Command Palette with \textbf{Ctrl+Shift+P} (Windows) or
\textbf{Cmd+Shift+P} (Mac). Search for \textbf{Create: New Jupyter Notebook}.
\item Save it as \texttt{setup\_check.ipynb} in \texttt{RTS-Research}.
\item Click \textbf{Select Kernel} at the top right, choose
\textbf{Python Environments}, and select the project's \texttt{.venv}.
\item Paste this into a \textbf{code cell} and press \textbf{Shift+Enter}:
\end{enumerate}
\begin{lstlisting}[language=Python]
import sys
import numpy as np
import pandas as pd
import scipy, seaborn, sklearn, openpyxl
import matplotlib.pyplot as plt
from sklearn.linear_model import LinearRegression

print("Python:", sys.version.split()[0])
print("Environment:", sys.executable)
model = LinearRegression().fit([[1], [2], [3]], [2, 4, 6])
print("Test prediction:", round(model.predict([[5]])[0], 1))
plt.plot([1, 2, 3], [2, 4, 6], marker="o")
plt.show()
\end{lstlisting}
\textbf{Done when:} you see \texttt{Test prediction: 10.0}, a line plot, and an
environment path containing \texttt{.venv}. Save your notebook. You are ready!

\section*{If something goes wrong}
\textbf{Windows cannot find pymanager or py:} Close and reopen VS Code and
PowerShell after installing Python Install Manager. If \texttt{py} still
points to an older installation, ask your leader to check the Python launcher.

\textbf{Cannot find .venv:} Open your \texttt{RTS-Research} folder in VS Code,
open a new terminal, and rerun step 4. If a command fails, fix that error
before continuing.

\textbf{Notebook says a module is missing:} Select the project's
\texttt{.venv} kernel again. If needed, rerun the library installation line.

\textbf{Still stuck after 10 minutes?} Send your group leader a screenshot
of the full error and the command you ran.

{\footnotesize Official help:
\href{https://docs.python.org/3/using/windows.html}{Python on Windows}
\quad|\quad \href{https://code.visualstudio.com/docs/datascience/jupyter-notebooks}{Jupyter in VS Code}.
Next time, open the same folder and notebook; you do not need to reinstall.}
\end{document}
